
% ============================================================
% Project : THE PHILOSOPHY OF VIABILITY
% Book    : Book I — Ontology of Viability
% Book ID : VS.B01
% Chapter : 08 — Comparative Philosophical Examination
% Section : 01 — Method and Criteria of Comparison
% Asset ID: VS.B01.P01.C08.S01.SS05.P010.TL01
% Type    : TL — Table
% File    : comparative_evaluation_criteria.tex
% Purpose : Present the distinct criteria used to evaluate
%           historical and contemporary ontological positions.
% ============================================================

\begin{table}[htbp]
\centering
\caption{Criteria for the Comparative Evaluation of Ontological Positions}
\label{tab:ch08-comparative-evaluation-criteria}

\small
\renewcommand{\arraystretch}{1.3}

\begin{tabular}{|p{0.19\textwidth}|p{0.34\textwidth}|p{0.35\textwidth}|}
\hline
\textbf{Criterion}
&
\textbf{Guiding Question}
&
\textbf{Failure or Limitation to Detect}
\\
\hline

Conceptual precision
&
Are the principal concepts and distinctions used consistently and with adequate clarity?
&
Equivocation, conceptual conflation, or distinctions without an identifiable argumentative function.
\\
\hline

Ontological adequacy
&
Does the account address the ontological question it claims to explain?
&
Substitution of epistemological, linguistic, or representational claims for ontological explanation.
\\
\hline

Argumentative validity
&
Do the conclusions follow from the premises, and are the premises adequately defended?
&
Circularity, hidden premises, invalid inference, or unsupported generalization.
\\
\hline

Explanatory reach
&
What range of problems can the account explain within its stated scope?
&
Overextension, unexplained cases, or claims of universality unsupported by the argument.
\\
\hline

Treatment of plurality and difference
&
Does the account explain the relevant forms of multiplicity without premature reduction?
&
Confusion between components, manifestations, relations, aggregates, and identities.
\\
\hline

Treatment of change and continuity
&
What changes, what remains, and under which transformation or respect?
&
Unjustified identification of identity, invariance, persistence, continuity, and stability.
\\
\hline

Scope discipline
&
Does the account distinguish established results from conditional claims and open questions?
&
Conversion of provisional constructions into settled conclusions or expansion beyond the evidence.
\\
\hline

Historical and interpretive responsibility
&
Is the position reconstructed from appropriate texts and contextual evidence?
&
Anachronism, misattribution, reliance on unsupported secondary summaries, or distortion through translation.
\\
\hline

\end{tabular}

\end{table}

% ============================================================
% END OF ASSET
% ============================================================
